\documentclass{article}
\usepackage[UTF8]{ctex}
\usepackage{amsmath,amssymb,amsfonts}
\usepackage{enumitem}
\usepackage{geometry}
\geometry{a4paper, margin=1in}
% === 错题本排版优化（不参与题目标识） ===
\setlength{\parindent}{0pt}
\setlength{\parskip}{0.35em plus 0.10em minus 0.08em}
\setlength{\emergencystretch}{3em}
\setlist{itemsep=0.25em, topsep=0.40em, parsep=0pt}
\allowdisplaybreaks[2]
% === 错题本排版优化结束 ===
\begin{document}
\section*{北京工业大学2024年数学分析真题}
\section*{8．北京工业大学2024年数学分析试卷}

\subsection{题目1}
一．用致密性定理（有界数列必有收敛子列）证明闭区间连续函数可以达到最大值．

\subsection{题目2}
二．\(z=(x, y)\) 是方程 \(x+y+z=\mathrm{e}^{z}\) 确定的隐函数，求 \(\frac{\partial^{2} z}{\partial x \partial y}\) ．

\subsection{题目3}
三．求曲线积分

\[
I=\int_{L} \frac{y \mathrm{~d} x-x \mathrm{~d} y}{x^{2}+y^{2}}
\]

其中 \(L\) 是 \(\frac{x^{2}}{4}+\frac{y^{2}}{9}=1\) 从 \((2,0)\) 逆时针到 \((-2,0)\) 的一段曲线。

\subsection{题目4}
四．用定积分求极限

\[
\lim _{n \rightarrow \infty} \sum_{k=1}^{n} \frac{n}{n^{2}+k} \sin \frac{k}{n}
\]

\subsection{题目5}
五．已知在 \(f(x, y)\) 在区域

\[
D=\{(x, y) \mid a \leq x<+\infty, c \leq y \leq d\}
\]

内连续，且无穷积分 \(\varphi(y)=\int_{a}^{+\infty} f(x, y) \mathrm{d} x\) 在 \([c, d]\) 上一致收敛，试证明 \(\varphi(y)\) 在 \([c, d]\) 上连续．

\subsection{题目6}
六．计算

\[
I=\int_{0}^{+\infty} \frac{\mathrm{e}^{-x} \sin a x}{x} \mathrm{~d} x
\]

\subsection{题目7}
七．证明无穷积分 \(\int_{1}^{+\infty} \frac{\sin x}{x^{p}} \cos \frac{1}{x} \mathrm{~d} x\) 在 \(p>1\) 时绝对收敛， \(0<p \leq 1\) 时条件收敛．

\subsection{题目8}
八．计算

\[
\oint_{C}\left(y^{2}-z^{2}\right) \mathrm{d} x+\left(z^{2}-x^{2}\right) \mathrm{d} y+\left(x^{2}-y^{2}\right) \mathrm{d} z,
\]

其中 \(C:\left\{\begin{array}{l}x^{2}+y^{2}+z^{2}=1 \\ x+y+z=1\end{array}\right.\) 的方向和 \(x+y+z=1\) 的法方向 \(\vec{n}=(1,1,1)\) 构成右手螺旋系．

\subsection{题目9}
九、计算

\[
I=\int_{0}^{\frac{1}{2}} \mathrm{~d} x \int_{x}^{1-x} \sqrt{y^{2}-x^{2}} \mathrm{~d} y
\]

\subsection{题目10}
十．已知 \(f(x)\) 在 \([0,1]\) 存在连续导函数，\(d_{n}=\int_{0}^{1} f(x) \mathrm{d} x-\frac{1}{n} \sum_{k=1}^{n} f\left(\frac{k-1}{n}\right)\) ，证明：

\[
\lim _{n \rightarrow \infty} n d_{n}=\frac{1}{2}[f(1)-f(0)]
\]
\end{document}
